\section{Contribuciones y Proyectos}

\cventry
{2025--Actualidad}
{Colaborador y Desarrollador}
{Identidad -- Proyecto MxOS}
{Open Source}
{}
{
\begin{itemize}
  \item Desarrollo de bibliotecas en \textbf{Rust} para la validación y generación de identificadores oficiales mexicanos:
  CURP, NSS y RFC.
  \item Diseño de soluciones eficientes, auditables y seguras para procesos de identidad digital.
  \item Enfoque en estandarización, software libre y soberanía tecnológica.
  \item Repositorio: \url{https://github.com/Mx-OS/identidad}
\end{itemize}
}

\cventry
{2025--Actualidad}
{Colaborador}
{MX-OS -- Sistema Operativo Linux Mexicano}
{Open Source}
{}
{
\begin{itemize}
  \item Participación en la iniciativa estratégica para el desarrollo de un sistema operativo Linux mexicano.
  \item Diseño del arte gráfico, logotipo y manual de identidad visual del proyecto.
  \item Colaboración activa con la comunidad y definición de lineamientos visuales.
\end{itemize}
}

\cventry
{2023--Actualidad}
{Autor y Desarrollador Principal}
{MEXA}
{Proyecto Personal / Open Source}
{}
{
\begin{itemize}
  \item Desarrollo de un validador y generador de campos para tramitología mexicana
  (CURP, RFC, NSS, CLABE, entre otros).
  \item Implementación de pruebas automatizadas y despliegue continuo.
  \item Proyecto disponible en:
  \url{https://github.com/fitorec/mexa-py} \\
  Página del proyecto: \url{https://fitorec.github.io/mexa-py/}
\end{itemize}
}

\cventry
{2018}
{Colaborador}
{Chamilo LMS}
{Open Source}
{}
{
\begin{itemize}
  \item Desarrollo de nuevas funcionalidades para la gestión de cursos, usuarios y herramientas de aprendizaje.
  \item Refactorización masiva del código, modernizando la arquitectura y adoptando mejores prácticas en PHP.
  \item Optimización de consultas a base de datos y mejora del rendimiento general del sistema.
  \item Corrección de bugs críticos y vulnerabilidades de seguridad.
  \item Aseguramiento de compatibilidad con versiones recientes de PHP y dependencias.
  \item Repositorio: \url{https://github.com/fitorec/chamilo-lms/commits?author=fitorec}
\end{itemize}
}

\cventry
{2015}
{Colaborador}
{Javalite}
{Open Source}
{}
{
\begin{itemize}
  \item Actualización de dependencias y herramientas de construcción (Maven, JUnit).
  \item Correcciones de mantenimiento en configuración, documentación y scripts de build.
  \item Ajustes para mejorar compatibilidad y estabilidad con nuevas versiones de librerías.
  \item Repositorio: \url{https://github.com/javalite/javalite/commits?author=fitorec}
\end{itemize}
}

\cventry
{2015}
{Colaborador}
{Joomla! CMS}
{Open Source}
{}
{
\begin{itemize}
  \item Corrección de bugs en componentes del core relacionados con sesiones, autenticación y validación de datos.
  \item Implementación de mejoras de seguridad y estabilidad en módulos críticos.
  \item Ajustes de compatibilidad con nuevas versiones de PHP y librerías externas.
  \item Refactorizaciones menores para mantener estándares de codificación.
\end{itemize}
}

\cventry
{2012}
{Colaborador}
{LibreCAD}
{Open Source}
{}
{
\begin{itemize}
  \item Contribuciones en C++ para mejorar el proceso de instalación en sistemas Linux.
  \item Ajustes y mejoras en la traducción al español mexicano.
  \item Repositorio: \url{https://github.com/LibreCAD/LibreCAD/commits?author=fitorec}
\end{itemize}
}

\cventry
{2011--2012}
{Colaborador}
{CakePHP Framework}
{Open Source}
{}
{
\begin{itemize}
  \item Correcciones y mejoras en la documentación y comentarios del código.
  \item Solución de bugs relacionados con validaciones y comportamiento de componentes.
  \item Refactorizaciones menores y ajustes en pruebas para mejorar la estabilidad del framework.
  \item Repositorio: \url{https://github.com/cakephp/cakephp/commits?author=fitorec}
\end{itemize}
}


\cventry
{2010--2022}
{Autor y Desarrollador Principal}
{DiaSQL}
{Proyecto Personal / Open Source}
{}
{
\begin{itemize}
  \item Plugins para el modelador de diagramas de base de datos.
  \item Permite exportar tus diagramas Dia a esquemas SQL.
  \item Proyecto disponible en:
  \url{https://github.com/fitorec/diasql}
\end{itemize}
}
